\section{Published operating points and the native noise control}
Antunes et al.\ train an autoregressive language model on crystallographic information files \cite{crystallm}. It generates a sequence whose decoded text specifies a lattice and sites. Wyckoff Transformer instead models the discrete occupation of crystallographic positions \cite{wyformer}. In our orbit experiment, the discrete occupation is an input obtained from a training-only catalog. The learned update acts on its continuous geometry. These differences determine what can be counted as a single generation step.
Dey et al.\ report 78.38\% match rate at 20 draws with one evaluation per candidate and 83.64\% with five. Their generation times for 10,000 structures on an A100 are 0.39 and 1.31 minutes, respectively \cite{uflow}. Our local replay of the released checkpoint recovers 79.32\% at one call and 20 draws. The different seeds, hardware and timing boundaries belong to the corresponding measurements. Table~3 of their paper compares 1,000-step DiffCSP and 100-step CrystalFlow at 20 draws, with 77.93\% and 78.34\% recovery respectively. Previously we also ran the released CrystalFlow checkpoint locally at five unannealed Euler steps, recovering 7,768/9,046 targets \cite{crystalflow}. For supplied Hall settings and occupied orbits, the separate DiffCSP++ replay recovered 121/128 at 20 draws of 40 evaluations each. These target labels provide information unavailable in formula-only prediction \cite{diffcspplus}.

The zero-call native uFlowCSP seed has nine independent Gaussian lattice entries. At 12 sites its observed mean cell volume was \(0.03158\,\text{\AA}^3\), compared with the predicted \(0.03169\,\text{\AA}^3\). Of 180,920 native draws, 141,900 had positive volume below \(0.1\,\text{\AA}^3\) and remained eligible for scale-enabled matching. Appendix~\ref{app:benchmark-conventions} derives this distribution. The symbolic prior starts from a training-calibrated volume, so the two learned updates face different initialization errors.
